\chapter[\emj{📦}~Arrays (Vectores)]{\emj{📦}~Arrays (Vectores)}

\begin{dsaquees}

Los vagones de un tren 🚂 numerados en orden. El vagón 0, el vagón 1,
el vagón 2... Cada vagón tiene exactamente UN pasajero (dato).

Para subir a cualquier vagón, solo dices el número (el índice) y 
¡listo! No tienes que pasar por todos los vagones anteriores. 
Esto es posible porque todos los vagones están pegaditos uno tras
otro en la vía.

Un array es una colección de elementos del mismo tipo guardados en 
posiciones de memoria CONTIGUAS (una tras otra).

\end{dsaquees}

\begin{lstlisting}
El Array es la estructura de datos más básica y fundamental. Su
característica principal es el acceso aleatorio (Random Access) en
tiempo constante O(1). 

En lenguajes modernos como C++, solemos usar "Dynamic Arrays" 
(como `std::vector`), que pueden crecer automáticamente si se llenan,
aunque por debajo siguen siendo un bloque de memoria contigua.
\end{lstlisting}

\begin{dsadiagrama}
\begin{verbatim}
Índice:  [0] [1] [2] [3] [4]
Valor:   [ 5][ 2][ 8][ 1][ 9]
          ↑
    Dirección de memoria contigua →→→→→→→→→→

Insertar 99 en el índice 2:
ANTES:   [ 5][ 2][ 8][ 1][ 9]
          ↓   ↓   ← desplazar →
DESPUÉS: [ 5][ 2][99][ 8][ 1][ 9]   ← O(n) porque movimos 8, 1 y 9
\end{verbatim}
\end{dsadiagrama}

\begin{lstlisting}
1. Acceso: Vas directo al índice `i`. Es instantáneo.
2. Búsqueda: Tienes que revisar uno por uno hasta encontrar el valor.
3. Inserción:
   - Al final: Si hay espacio, es instantáneo.
   - Al inicio/medio: Tienes que "empujar" a todos los de la derecha
     un lugar para hacer espacio.
4. Eliminación:
   - Al final: Borras y ya.
   - Al inicio/medio: Tienes que "jalar" a todos los de la derecha
     un lugar a la izquierda para cerrar el hueco.

PSEUDOCÓDIGO
──────────────────────────────────────────────────────────────

// Insertar en medio
función insertar(array, valor, posicion):
    PARA i desde tamaño hasta posicion + 1:
        array[i] = array[i-1] // desplazar a la derecha
    array[posicion] = valor

// Eliminar en medio
función eliminar(array, posicion):
    PARA i desde posicion hasta tamaño - 2:
        array[i] = array[i+1] // desplazar a la izquierda

CÓDIGO C++
──────────────────────────────────────────────────────────────

#include <iostream>
#include <vector>   // En C++ usamos vector para arrays dinámicos
#include <algorithm>
using namespace std;

int main() {
    // 1. Declaración e inicialización
    vector<int> arr = {5, 2, 8, 1, 9};

    // 2. Acceso por índice -> O(1)
    cout << "Elemento en idx 2: " << arr[2] << endl; // 8

    // 3. Insertar al final -> O(1) amortizado
    arr.push_back(10); // [5, 2, 8, 1, 9, 10]

    // 4. Insertar en medio -> O(n)
    // Insertamos el 99 en la posición 2
    arr.insert(arr.begin() + 2, 99); // [5, 2, 99, 8, 1, 9, 10]

    // 5. Eliminar en medio -> O(n)
    // Borramos el elemento en la posición 4
    arr.erase(arr.begin() + 4); // [5, 2, 99, 8, 9, 10]

    // 6. Tamaño del array
    cout << "Tamaño: " << arr.size() << endl;

    // 7. Recorrer el array -> O(n)
    for(int x : arr) {
        cout << x << " ";
    }
    cout << endl;

    return 0;
}

COMPLEJIDAD (Big-O)
──────────────────────────────────────────────────────────────

  Operación        │ Tiempo    │ Espacio
  ─────────────────┼───────────┼─────────
  Access           │ O(1)      │ O(1)
  Search unsorted  │ O(n)      │ O(1)
  Search sorted    │ O(log n)  │ O(1)
  Insert end       │ O(1)*     │ O(1)
  Insert middle    │ O(n)      │ O(1)
  Delete end       │ O(1)      │ O(1)
  Delete middle    │ O(n)      │ O(1)
  *amortizado en dynamic array (vector)
\end{lstlisting}

\begin{dsaentrevistas}

✅ Arrays son la estructura de datos más común en entrevistas. Domínalos.

💡 Patrones clave:
   - "Two Pointers": Usar un índice al inicio y otro al final.
   - "Sliding Window": Una ventana que se desliza para encontrar subarrays.
   - "Prefix Sum": Pre-calcular sumas para responder rangos en O(1).

⚠️ Off-by-one errors: El error más común es salirte del rango (acceder
   a arr[n] cuando el tamaño es n). Siempre usa `i < n`.

🔑 Pregunta clave: "¿El array está ordenado?". Si está ordenado, casi
   siempre la respuesta involucra Binary Search O(log n).

\end{dsaentrevistas}

\begin{lstlisting}
• Array = bloque de memoria contigua.
• Acceso instantáneo O(1) si conoces el índice.
• Insertar/Eliminar en medio es LENTO O(n) por los desplazamientos.
• En C++: `std::vector` es el rey.
• Muy eficiente en espacio (mínimo overhead).
• Ideal para cuando conoces el tamaño o solo agregas/quitas al final.
════════════════════════════════════════════════════════════════
\end{lstlisting}


